\subsection{新型透射电镜}

前面几节的透射电镜回答的是「催化剂是什么」；本节要回答的是「在什么条件下还能看到什么」，以及每一项新能力把限制推到了哪里。全节的组织原则只有一条：每一项新技术都写成一对——它解除了什么旧限制，又引入了什么新限制。前者决定它值得用在什么地方，后者决定它的结论可以被推到多远。读懂这一对，比记住任何一台电镜的分辨率指标都重要。

\subsubsection{球差校正的透射电镜与扫描透射电镜}

高分辨像（晶格条纹像）是现代透射电镜中最常用的技术，它能够提供微米或纳米尺度催化剂的原子结构细节；但物镜球差会使这些结构细节变模糊。球差是由电磁透镜中心区域和边缘区域对电子会聚能力不同而造成的：远轴电子通过透镜时被折射得比近轴电子厉害得多，因而由同一物点散射的电子经过物镜后不能汇聚在一点上，而是在物镜像平面上变成了一个漫散圆斑（球差系数 $C_s>0$）。球差降低成像质量表现在两个方面：（1）图像中漫散点相互重叠，导致图像模糊；（2）两束不同的电子波相互作用，在图像的边缘区域扩展形成菲涅耳条纹，使试样的真实表面不清。

在球差校正器商业化之前，通行的绕行方案是利用一系列不同欠焦量的图像，通过计算重建样品的真实结构（出射波重构）。这条路可行，但这是相当费时间的工作：它把像差问题转化成了「大量采图 $+$ 大量计算」的问题，而且重建质量依赖焦系列的间隔与对中精度。

\begin{principle}
球差在像面上造成的是一个随会聚半角 $\alpha$ 三次方增长的漫散半径
\begin{equation}
\delta_s \approx C_s\alpha^{3}
\label{eq:ch02c-cs}
\end{equation}
式~\eqref{eq:ch02c-cs} 说明旧技术面临的是一个权衡：要压小球差模糊就得缩小会聚角，而缩小会聚角又使衍射斑变大（衍射极限与 $\lambda/\alpha$ 同阶），两者不能同时改善。球差校正的价值正是在于解除这个权衡——不必再靠小光阑换清晰度。

现在商业化的校正器有两类，分别对准成像与束斑。成像校正器（imaging Cs corrector）放在物镜的像平面之后，校正后的像由中间镜继续放大；注意这句话的推论：它校正的是物镜三阶球差对像的影响，并不取消校正器之后的透镜。束斑校正器（probe corrector）面向扫描透射模式：要用 STEM 得到原子像和原子的电子能量损失谱，电子束必须会聚为一个原子大小的电子探针，同时还要保证探针有一定的强度，以避免无限制地延长记录透射信号的时间；但由于电磁透镜存在像差，不可能既把入射电子束形成最小束斑又保证其有一定强度。

束斑校正器利用四极或六极、八极线圈产生负的球差来抵消透镜组的球差。据源文对四极—八极联用光路（图~2-49）的说明：聚光镜产生的束斑带有正向球差，第一级四极磁透镜 Q1 先把它在 $x$ 方向拉伸成高度椭圆形截面，再由八极磁透镜 O1 在 $y$ 方向叠加散焦作用（负球差），并由第二级四极磁透镜 Q2 把束斑复形；于是在第二级八极磁透镜 O2 的平面上，得到了与原聚光镜出射束斑形状相同、但在 $y$ 方向叠加了负球差的光斑，完成了 $y$ 方向的球差调整。同理，在后续光路中运用第三级四极、八极及第四级磁透镜的组合，完成光斑在 $y$ 方向的拉伸和 $x$ 方向的散焦作用。最终在物镜前焦面处获得与原来形状相同、但对离轴光线叠加了散焦作用的束斑，使整个束斑内的光线全部能汇聚一点。
\end{principle}

此处已对照原页面图像（书内 p.147 / PDF p.166）：原书只写球差使同一物点散射的电子在物镜像平面上「变成了一个漫散圆斑（球差系数 $C_s$ 大于零）」（图~2-48），\textbf{没有给出漫散半径的公式}，也没有出现四阶项或波前像差记号；式~\eqref{eq:ch02c-cs} 是本改编本按三阶几何像差的标准形式补出的（$\delta_s\approx C_s\alpha^{3}$），不是原书内容。

\begin{measure}
球差校正带来的性能提升，源文给出的量是：应用束斑校正器并通过计算机自动调整后，HAADF-STEM 的 $Z$ 衬度分辨率可从 0.24\,nm 提高到 0.13\,nm，同时电子束斑强度增加了一个数量级，从而降低噪声背底；如果能够进一步修正消除色差和更高阶的几何球差，束斑直径可望提高到 0.05\,nm。成像校正器一侧，源文只给出「分辨率优于 0.1\,nm」这一量级，并强调校正后可以得到没有透镜畸变影响的高分辨电子显微图像。

这几项指标都必须与记录条件捆绑使用。本节归纳的记录与报告项为：加速电压、球差校正器的设置（校正到哪一阶）、欠焦量、会聚角或探针会聚半角、探针电流与像素驻留时间、总电子剂量、样品厚度估计，以及用于标定分辨率的标样（如金颗粒、石墨烯）。只报「分辨率 0.1\,nm」而不报这些条件，读者无法判断该指标是否适用于自己的样品。
\end{measure}

\begin{dataflow}
校正后的像 $\to$ 结论的链条是：原始图像（含扫描畸变与非均匀性）$\to$ 漂移与畸变校正 $\to$ 由标样标定放大倍数 $\to$ 判读原子柱位置或晶格间距 $\to$ 结合衍射确认物相 $\to$ 结合 HAADF 强度趋势或 EDS/EELS 确认元素。这条链条上任何一环缺失，原子柱位置的判读都只能算「与某结构模型一致」，不能算「测出了该结构」。
\end{dataflow}

\begin{pitfall}
最常见的误判是把球差校正后的分辨率当成「能看到任何原子」。点分辨率是一个在特定样品、特定条件下可分辨最小间距的指标，它不等于「任意原子都能看见」。至少在以下四种情况下它不成立：

（1）轻元素与低 $Z$ 载体。HAADF 的衬度只是原子序数的单调函数（见 \S 2.7.3 的式~\eqref{eq:ch02c-haadf}），氧、碳这类轻原子柱在重原子柱旁边或载体之上时，衬度差可以小于噪声；「看到」的前提是衬度差大于噪声，而不是仪器分辨率足够小。

（2）样品厚度。动力学衍射与多重散射在厚区使强度不再随原子序数、厚度单调变化，投影势的解释失效；高分辨像中的暗（或明）条纹不是试样某一晶面的原子排列，而只是具有这一晶面间距的条纹像。

（3）束敏样品。要达到原子分辨所需的剂量往往已经超过样品的耐受剂量，图像上的「结构」可能是电子束诱导的重构、非晶化或物种迁移的产物（详见 \S 2.8 关于束损伤的讨论）。

（4）校正不等于无像差。球差校正器消除的是三阶球差；色差与更高阶几何像差仍是独立的分辨率瓶颈，源文也正是把「消除色差和更高阶几何球差」作为进一步提高束斑直径的前提。
\end{pitfall}

\begin{crosscheck}
球差校正像给出的结论应当用三条独立证据回查：一是电子衍射（衍射来自大量粒子，物相判定比单张像更具代表性）；二是 HAADF 强度趋势或 EDS/EELS 的元素与价态信息（把「有原子柱」升级为「是什么元素的原子柱」）；三是同一区域在不同欠焦量或不同剂量下的重复成像（区分真实结构与束诱导结构）。三者一致时，本节所述的原子尺度结论才站得住。
\end{crosscheck}

\subsubsection{高能量分辨率扫描透射电镜：单色器与高分辨能量过滤器}

配备了电子能量损失谱仪的透射电镜，在化学分析和结构判定上比单纯使用电子衍射和高分辨像有更高的灵活性。用 ELNES 能够在纳米尺度内分析各种材料的局域化学成分和电子结构，但是热场发射 TEM 系统的能量分辨率通常只有 0.8--1\,eV，与电子结构相关的精细结构会湮没在谱中。如果能提高谱的分辨率并与理论计算比较，从电离损失峰近阈精细结构得到的信息的可靠性将大大提高——这里要解决的是一个「能量」限制，而不是空间限制。

提高能量分辨率主要有以下几种方法：一是在肖特基场发射枪的前面放置 Wien 型能量单色器，降低电子能量的分散度；二是通过电子衰减装置提高高压箱的稳定性；三是采用高分辨能量过滤器（如在磁棱镜前面添加多极线圈以消除第三级和第四级谱仪像差的商用过滤器），这类过滤器不仅能减小零损失峰位移，而且可以提高电子光学性能。装备了 Wien 单色器、稳定的高压箱和高分辨能量过滤器的透射电镜，能量分辨率可高于 0.1\,eV。源文对这一进展的评价是：这将开辟高分辨电子能量损失谱在固体物理、固体化学和材料科学领域应用的新纪元。

\begin{principle}
单色器提高的不只是谱的能量分辨率，它还通过时间相干性改变成像的信息极限。当入射电子的能量分布变窄后，电子束的时间相干性提高，图像的衬度随之提高；在衬度传递函数中，表现为包络函数变宽且下降变缓，倒空间中的频率传播增大，于是透射电镜的信息分辨极限被扩展。源文把这一包络函数记为式（2-13）中的 $D$，本节按电子显微学的标准形式把它写成
\begin{equation}
E_c(u)=\exp\!\left[-\left(\frac{\pi\lambda\Delta u^{2}}{2}\right)^{2}\right],
\qquad
\Delta \approx C_c\frac{\Delta E}{E}
\label{eq:ch02c-envelope}
\end{equation}
其中 $u$ 为空间频率，$\lambda$ 为电子波长，$\Delta$ 为欠焦（聚焦）扩展，$C_c$ 为色差系数，$\Delta E/E$ 为入射电子能量的相对分散。令指数等于 2 可得信息极限 $u_{\max}\propto(\lambda\Delta)^{-1/2}$，即 $d_{\min}=u_{\max}^{-1}\propto(\lambda\Delta)^{1/2}$：单色器把 $\Delta E$ 减小一个数量级，时间相干包络单独作用下的信息极限约改善 $\sqrt{10}\approx 3$ 倍。
\end{principle}

此处已对照原页面图像（书内 p.105 / PDF p.124 与书内 p.151 / PDF p.170）：原书式（2-13）印为 $A(u,v)=\exp\{2\pi\mathrm{i}\chi(u,v)\}B(u,v)D(u,v)$，式中 $B(u,v)$ 为照明电子束发散度引起的衰减包络函数，$D(u,v)$ 为色差效应引起的衰减包络函数；\textbf{原书只给出 $B$、$D$ 的物理含义，没有写出二者的显式表达式}（紧随其后的式（2-14）给出的是相位差 $\chi(u,v)$ 的显式形式）。故本节式~\eqref{eq:ch02c-envelope} 的显式形式是按电子显微学的标准「时间相干包络」补出的，符号 $D$ 沿用原书式（2-13）的记号。原书在 §2.7.2 只说明包络函数「将变宽且下降变缓，倒空间中的频率传播增大」，以及使用单色器的代价——「能量分辨率小于 0.1\,eV 时，将有 75\% 的场发射电子被单色器滤除掉」。

\begin{example}[五氧化二钒中钒 L3 边的 ELNES 与 NEXAFS 对照]
源文用一台装备了 Wien 单色器、稳定高压箱和图像过滤器的透射电镜采集零损失峰，在较长采谱时间下能量分辨率仍达 0.22\,eV（图~2-50）。随后用 0.22\,eV 分辨率采集五氧化二钒晶体中钒的 $L_3$ 电离峰的近阈精细结构（ELNES），并与用 0.08\,eV 能量分辨率采集的钒 $L_3$ 的 X 射线吸收近阈精细结构（NEXAFS）作对照（图~2-51）。下面把「图上直接读到的量」与「据此作出的解释」严格分开。

\begin{obs}
（1）0.22\,eV 的 ELNES 与 0.08\,eV 的 NEXAFS 在钒 $L_3$ 电离边处都出现一列标记为 P1--P6 的峰。

（2）P1 与 P2 两峰的间距：NEXAFS 为 1.1\,eV，ELNES 为 1.2\,eV。

（3）NEXAFS 在 1.5\,eV 处有 P6 峰，而 ELNES 谱中该位置的峰非常弱。

（4）当不使用 Wien 单色器时（能量分辨率约 1\,eV），钒 $L_3$ 上的 P1--P6 精细结构不再可分。

（5）两组谱的峰位与相对强度并不完全重合，ELNES 与 NEXAFS 的谱形存在可见差异。
\end{obs}

\begin{infer}
（1）P1、P2 峰由钒 $3d$ 导带的劈裂引起；实验得到的劈裂间距 1.1\,eV（NEXAFS）与 1.2\,eV（ELNES）和量子力学计算的约 1\,eV 符合得很好。

（2）P4、P5 可以由钒 $3d$ 中平台状分布的 $e_g$ 态（集中在 4.5--7\,eV，其中 5\,eV 与 6.5\,eV 两个特征）解释，尽管实验谱结构与模型的精细结构并不完全相同。

（3）NEXAFS 谱中位于 1.5\,eV 处的 P6 峰（在 ELNES 中非常弱）不能用钒 $3d$ 导带的未占据态解释，它可能是目前能带计算中没有考虑的表面态引起的；P3 峰的来源，现在的能谱计算还不能解释。

（4）最有信息量的推论与「分辨率的直觉」相反：NEXAFS 的能量分辨率（0.08\,eV）远高于 ELNES（0.22\,eV），但并没有揭示更多的精细结构峰特征。也就是说，把仪器分辨率再提高，并不自动换回更多可分的信息。源文没有进一步解释这一现象；电子显微学的通行解释是，过渡金属 $L$ 边的本征寿命展宽（核空穴寿命展宽）量级约 0.1--0.3\,eV，当仪器分辨率优于该尺度后，谱线宽度由样品本身决定，继续提高分辨率不再增加可分峰数。因此读谱时的正确顺序是：先判断限制来自仪器还是样品本身，再决定要不要为提高分辨率付出代价。
\end{infer}
\end{example}

提高能量分辨率后的另一个新应用是用高分辨能损谱研究纳米半导体材料的能隙；但最近的研究表明，相对论效应和电子的切伦科夫辐射会导致用能谱测量能隙宽度出现很大的偏差。源文由此给出的方法论判断是：当使用一系列综合技术提高能量损失谱的分辨率后，透射电镜中配备的能损谱能够提供与同步辐射技术相同的信息，可以把对表面敏感的同步辐射方法和对体相敏感的电子能量损失方法结合起来研究材料的电子结构，同时还能利用电子衍射和高分辨像得到纳米级分辨率的完整结构信息——这同时也构成了对电子能带计算理论与方法的挑战。

\begin{measure}
能量分辨率的实际可用性由三个量共同决定：仪器能量分辨率、可用的入射束流、以及采集一张谱所需的时间。单色器是以非常小的入射电子束强度为代价的：源文明确给出，当能量分辨率优于 0.1\,eV 时，将有约 75\% 的场发射电子被单色器滤除掉。束流下降到四分之一，意味着在相同的信噪比下采集时间约为原来的四倍，而总剂量是束流密度与时间的乘积，因此束敏样品在这种模式下受到的总剂量并不因为「分辨率更高」而下降。报告高分辨 EELS/ELNES 结果时，至少应给出：能量分辨率及其标定方式（零损失峰半高宽）、色散（每通道 eV 数）、采集时间与束流、样品厚度参数 $t/\lambda$、以及是否做过背底扣除与多重散射去卷积。
\end{measure}

\begin{dataflow}
从原始谱到电子结构结论的链条为：原始低能损失与高能损失谱 $\to$ 零损失峰对齐与能量标定 $\to$ 背底扣除（近阈区常用幂律或指数形式）$\to$ 去除多重散射（傅里叶对数去卷积等）$\to$ 与参考化合物谱、或与未占据态密度计算对照 $\to$ 给出价态、配位对称性与未占据态特征。链条上每一步都带入误差，而「与计算符合得很好」（如本节 1.1--1.2\,eV 对 1\,eV）只是相容性证据，不是唯一性证据：不同的模型细节都可能给出相近的峰间距。
\end{dataflow}

\begin{pitfall}
（1）把「分辨率更高」直接读成「信息更多」。上面的实例已经证明两者不等价：0.08\,eV 的 NEXAFS 没有比 0.22\,eV 的 ELNES 多给出峰特征，因为此时限制改由样品本身的本征展宽决定。

（2）忽略束流代价。单色器滤掉约 75\% 的场发射电子，是本项技术的固有代价；在束敏样品上，为了分辨率而延长的采集时间会直接转化为束损伤。

（3）在低能区做定量而不考虑相对论与切伦科夫效应。源文明确指出，用能谱测量纳米半导体能隙宽度时，相对论效应和电子的切伦科夫辐射会导致很大的偏差，因此基于低能损失谱的带隙数值必须做相应修正后才可比较。
\end{pitfall}

\begin{crosscheck}
高分辨 EELS 的结论应与下列证据交叉：同一元素的 X 射线吸收谱（同步辐射或 NEXAFS，对表面敏感而 EELS 对体相敏感）、电子衍射确定的物相与取向、以及不同加速电压下的重复测量（改变切伦科夫辐射与多重散射条件）。两组谱一致时，价态与未占据态的解释才可信；不一致时，首先要检查的是厚度、去卷积与能量标定，而不是急着修改模型。
\end{crosscheck}

\subsubsection{三维电子显微术：电子断层重构}

TEM 和 HRTEM 有一个共同的缺点：只能对三维物体进行二维投影成像。对于一些材料，二维投影分析已经足够；但是对于许多催化剂，二维投影就是一个局限。许多催化剂通常具有非常复杂的结构，如表面粗糙、微孔介孔、规则和不规则的通道结构，二维投影会把这些结构特征相互覆盖，使获得三维结构信息变得困难。源文记载，首次对催化剂进行三维表征的是荷兰 Utrecht 的 Bram Koster 等；这一方法起源于 20 世纪 60 年代的 X 射线 CT，后来推广到电子显微镜上，但主要用来分析生物分子的三维结构，成为具有纳米级分辨率的分析大分子的一个常用工具，近年才被用于材料科学但已得到相应的普及。

三维电子显微术（电子断层扫描，electron tomography）的基本过程就是转动样品，每转动一个角度就记录一张电子显微像，之后用这一系列照片（通常至少要一百张或更多）做重构计算，得出样品的三维结构图。

\begin{principle}
断层重构的数学基础是投影与反投影。以 $f(x,y)$ 记样品的某个与成像强度成正比的物理量（在 HAADF 模式下近似为质量厚度），角度 $\theta$ 下的一维投影写成
\begin{equation}
p_\theta(t)=\iint f(x,y)\,\delta\!\left(x\cos\theta+y\sin\theta-t\right)\mathrm{d}x\,\mathrm{d}y
\label{eq:ch02c-radon}
\end{equation}
投影切片定理指出，一维投影的傅里叶变换等于物体二维傅里叶变换中过原点、与投影方向垂直的一条切片；因此只要倾角覆盖足够完整，就可以由投影序列反演出 $f(x,y)$。

要让式~\eqref{eq:ch02c-radon} 的前向模型成立，每张电子显微图的强度必须是试样在该角度下质量（质厚）的单调函数，也就是说显微像中不能含有衍射衬度——衍射衬度只在某些满足布拉格条件的角度出现，其强度不是试样质厚的线性函数。源文给出的解决办法是用 HAADF 高角环状检测器收集散射角大于布拉格衍射角的高角度散射电子成像，这样可以排除衍射衬度，从而可近似认为像的强度与试样的厚度和组成元素的原子序数平方成正比
\begin{equation}
I_{\mathrm{HAADF}}\ \propto\ Z^{\alpha}\,t,\qquad \alpha\approx 2
\label{eq:ch02c-haadf}
\end{equation}
其中 $t$ 为局部厚度。图~2-52 给出了在 STEM 上做三维电子成像的基本原理示意：用透射电子束扫描样品，由 HAADF 检测器搜集从不同高角度透射过样品的电子成像。
\end{principle}

此处已对照原页面图像（书内 p.152 / PDF p.171）：原书 §2.7.3 只用文字交代过程——「三维电子显微技术（electron tomography）的基本过程就是转动样品，每转动一角度就记录一张电子显微照片。之后用这一系列照片（通常至少要一百张或更多）做重构计算得出样品的三维结构图」——\textbf{既没有给出投影（Radon 变换）的积分式，也没有提到投影切片定理}；式~\eqref{eq:ch02c-radon} 与投影切片定理一句是按断层重构的标准形式补写的，记号 $f$、$p_\theta$ 与归一化为本改编本约定。

关于 HAADF 条件，原书原话是「通过 HAADF 高角环状检测器搜集起来的、散射角大于布拉格衍射角的高角度散射电子像可以排除衍射衬度；这样的话，可近似认为像的强度与试样的厚度和组成元素的原子序数平方成正比」。原书直接写\textbf{「原子序数平方」}，未引入指数符号，也未讨论指数随收集角变化；式~\eqref{eq:ch02c-haadf} 把这一表述记作 $Z^{\alpha}$ 并写 $\alpha\approx 2$，而「实际 $Z$ 指数通常介于 1.7 与 2 之间并随收集角变化」是本改编本补出的学科注记，不是原书内容。同页图~2-52 的图题为「电子投影技术的基本原理示意图」。

\begin{measure}
三维重构对硬件提出了两项互相牵连的要求。

（1）样品台必须能大角度倾转。源文指出，照三维电子显微照片需要用特殊制造的样品台，以保证试样能正负转至少 80°，而这要加大物镜极靴的空间以便样品转动，这会导致物镜的分辨率下降；通常的做法是加上前述的球差校正器提高分辨率，把这一损失弥补回来。这是一个典型的「新限制被另一项新技术补掉」的例子，也是本节把各项技术成对书写的原因。

（2）倾转序列本身必须稳定。原子分辨的三维重构依赖于整个样品倾转过程中的高度稳定性，以确保每张图片的分辨率；同时由于高倍成像下样品承受高电子辐照剂量，重构往往采用低加速电压结合 probe 校正器的搭配。源文给出的典型工作是美国劳伦斯伯克利国家实验室在 80\,kV 下获得铝基底中内嵌银单晶纳米颗粒的原子分辨重构图像。值得一提的是，受限于材料对电子辐照耐受度，目前还无法在确保原子分辨条件下实现元素乃至价态区分，而只能实现单一元素的原子结构重构。

重构实验至少应报告：倾角范围（如 $\pm 80^\circ$）、投影张数、对位方法、是否做了强度归一化、重构算法（反投影、SIRT 或迭代类方法）、体素分辨率判据、以及分割阈值如何选取。
\end{measure}

\begin{dataflow}
倾转序列 $\to$ 三维结论的链条为：采集投影序列（$\geq 100$ 张）$\to$ 对位（互相关或标记点）$\to$ 强度归一化（补偿厚度与检测器响应）$\to$ 反投影或迭代重构 $\to$ 分割与定量（孔内/孔外金属粒子比例、粒径与空间分布、核壳结构的价态分区）$\to$ 报告不确定度。定量分析可以得到金属粒子在载体表面和在孔径内分布的比例，这是二维投影做不到的。
\end{dataflow}

源文把三维重构的发展归纳为两个方向：更高的空间分辨率，即原子分辨三维重构；更高的能量分辨率，即元素／价态分辨三维重构。后者主要体现在两条路上：一是基于 EELS 的等离子基元三维重构（低能光学吸收行为）及元素价态的三维重构（元素吸收边的精细结构）；二是基于高效 EDS 的多元素组分三维重构。在 EDS 一侧，目前的趋势是采用对称分布的多个 EDS 探头，例如配置 4 个 SDD 探头，可显著增加信号接收立体角，提升能谱探测效率，缩短重构时间从而降低样品的辐照损伤概率；对称分布的探头还可以有效减弱样品倾转过程中周边遮挡引发的阴影效应，提升元素定量精度。

三维电子显微术在催化中最重要的应用是测定催化剂的三维结构、纳米粒子的聚合、催化剂活性位（金属粒子）在孔状载体的位置和空间分布。源文举出的实例包括：基于带边精细 EELS 谱重构手段获得的 30\,nm FeO/Fe$_3$O$_4$ 核壳结构中 Fe$^{2+}$ 和 Fe$^{3+}$ 的分布——对 Fe、Co 这类变价组分，其空间分布对准确把握反应中气体分子的吸附及活化至关重要；以及基于电子能量单色器的低能 EELS 三维重构，可以在小于 0.1\,eV 的尺度上获得单个粒子不同等离子基元模式（对应不同响应光波长）的空间立体分布，例如 100\,nm 银立方体的表面等离子基元分布，对应基元能量分别为 2.2、2.7、2.9、3.3 及 3.6\,eV。

\begin{obs}
在上述银立方体的低能 EELS 三维重构结果中，可以读出：（1）在 2.2--3.6\,eV 范围内存在五个可分辨的等离子基元模式；（2）这五个模式在立方体表面的空间分布并不相同；（3）立方体的四个顶角中有三个（源文表述为「上面 4 个顶角」之外的另一个顶角）的模式能量与其余顶角不同。
\end{obs}

\begin{infer}
该顶角模式能量与其余顶角不同，被解释为存在金属与基底的耦合作用。请注意这条推断的模型依赖：它是「重构出的模式空间分布 $+$ 等离基元耦合的电磁模型」共同给出的结论；重构本身给出的只是投影序列反演出的强度／能量分布，而不是一张可以直接量出耦合强度的「照片」。
\end{infer}

\begin{pitfall}
三维重构最容易被误读成「无模型的直接测量」。实际上，一条重构结果至少捆着四个假设：一是前向模型（式~\eqref{eq:ch02c-haadf} 的 $Z^\alpha t$ 单调性，要求无衍射衬度、样品厚度适中且检测器响应已知）；二是角度采样的完整性；三是反演算法的正则化与迭代终止条件；四是分割阈值。前两条带来的偏差具有系统性和方向性：

（1）缺失楔形。既然只能倾转 $\pm 80^\circ$，沿电子束光轴方向就缺失了一段角度信息，重构在该方向上的分辨率与形貌会发生失真（孔道、颗粒沿光轴被拉长或用伪影填补）。报告三维结果时若不给倾角范围与缺失楔形，读者无法判断某个「通道」或「空腔」是真实的还是反演伪影。

（2）投影之间的不一致。倾转过程中的样品漂移、局部厚度变化、衍射衬度残留、EDS 的阴影效应与吸收效应，都会被反演算法当成「结构」写进结果里。

（3）能力边界。原子分辨与元素／价态分辨目前不能同时满足：能给出原子分辨重构的场合，往往只能对单一元素成像；能给出元素／价态分区（如 Fe$^{2+}$/Fe$^{3+}$）的场合，空间分辨率要低得多。

因此三维重构给出的「位置」和「比例」是带模型误差与角度偏差的估计量，而不是直接测量值。凡是把重构结果当作无模型直接测量、跳过上述四个假设去外推孔径分布或活性位密度的结论，都不成立。
\end{pitfall}

\begin{crosscheck}
三维重构结果应当用三类证据回查：一是同一体系的二维投影像（沿特定晶带轴的 HRTEM 可以验证重构中同一位置的晶格与界面）；二是元素／价态信息的另一条独立通道，即 EDS 与 EELS 对同一区域的一致性（两者对厚区的响应不同，一个受吸收与荧光影响，一个受多重散射影响）；三是物理合理性检查，例如由重构得到的孔内／孔外粒子比例是否与化学吸附测得的金属分散度、孔容与孔径分布自洽。三条中若有两条冲突，先怀疑重构的倾角完整性与分割阈值，而不是先改结论。
\end{crosscheck}

\subsubsection{电子全息成像}

除了校正球差能提高电子显微镜的点分辨率外，利用电子全息术也能提高透射电子显微镜的分辨率。电子全息术的基本原理是 Gabor 在 1948 年为试图超越电子显微镜点分辨率时提出的：以电子波的相位为基础产生没有任何透镜畸变的真实图像，包括振幅像和相位像。一般的高分辨电子显微镜中出射的电子波的相位与振幅像进行了非线性混合，且最终的像衬度受物镜相位传递函数的复杂修饰；电子全息术则是通过恢复像的相位和振幅，进而再生真实物体信息的技术。

实验由两个阶段组成：一是电子显微镜记录电子的相位信息，即在电镜中装入双棱镜，拍摄电子波的相干条纹（全息照片）；二是它的再生，即对全息图像进行傅里叶变换而再生像的相位与振幅。电子全息术的实验基础是高亮度、高相干性的电子源，这正是现在场发射电子显微镜所具备的，因此最近十几年来电子显微镜中的电子全息术推广得很快。

\begin{principle}
对薄样品，出射电子波的相位与样品内的投影电势近似成线性关系
\begin{equation}
\varphi(x,y)=C_E\int V_0(x,y,z)\,\mathrm{d}z
\label{eq:ch02c-holo}
\end{equation}
其中 $V_0$ 为晶体的平均内电势，$C_E$ 为与加速电压和电子波长有关的相互作用常数。式~\eqref{eq:ch02c-holo} 就是全息术能够「测量场」的原因：相位是投影电势的线积分，因此从重建出的相位图可以反推厚度、电势与场分布。除消除物镜球差的影响外，源文列出三类应用：研究带电粒子周围的静电场分布、纳米粒子的投影电势像的厚度（后者可用来测定粒子是否包含硬的核心或空洞）、以及研究磁畴的结构和其周围的场分布。
\end{principle}

此处已对照原页面图像（书内 p.154 / PDF p.173）：原书 \S 2.7.4 全节只作文字叙述——「电子全息术以波的干涉和衍射性质为基础产生没有任何透镜畸变的真实图像，包括振幅像和相位像」，随后只列出三类应用（带电粒子周围的静电场分布、纳米粒子的投影电势像的厚度、磁畴的结构和其周围的场分布），\textbf{并未给出相位与投影电势之间的关系式}，全节没有出现符号 $C_E$，也没有积分式。故式~\eqref{eq:ch02c-holo} 及其中的 $C_E$ 是本改编本按电子全息的标准形式补出的，非原书内容。另：原书该段把再生的对象印作「相位与辐射」，据上下文（前半句为「振幅像和相位像」）当为「振幅」之误，译文按「振幅」处理。

\begin{measure}
全息实验的记录条件包括：双棱镜偏压与位置、全息条纹间距与全息图视场、参考全息图的采集（用于扣除几何相位与透镜畸变）、以及重建所采用的边带或相位解缠方法。由于式~\eqref{eq:ch02c-holo} 依赖薄样品与弱相位物体近似，报告相位定量结果时还应给出样品厚度估计与其不确定度。在催化方面，源文明确说明这类应用还很有限，最近有文献报道用电子全息成像术研究 Au/TiO$_2$ 催化剂中金与二氧化钛界面处的电荷传递，试图解释金粒子低温对一氧化碳氧化的催化机理。
\end{measure}

\begin{pitfall}
（1）把「相位像」当成与结构一一对应的图。相位是投影电势（或磁场）沿束方向的积分，不同深度上的电荷与场会叠加成同一个相位值；只有在厚度、成分已知且磁贡献与静电贡献可分的前提下，相位才能换算成某一具体的电荷密度或场强。

（2）忽略相位解缠与视场限制。相位是周期性量，超过 $2\pi$ 的相位跳变需要解缠，解缠算法在低信噪比或强场区域会出错；同时双棱镜全息要求足够的相干宽度，视场越大对相干性的要求越高，这限制了能同时观测的区域尺寸。

（3）把式~\eqref{eq:ch02c-holo} 用在厚样品或强衍射条件下。动力学衍射使相位不再与投影电势成简单线性关系，此时由相位反推的厚度与电势会同时偏离。

因此电子全息在催化中的定位应当是「场与电荷传递的专用探针」，而不是又一种通用成像模式；源文对其在催化中的应用仍持谨慎态度，这一点应当被如实转述。
\end{pitfall}

\begin{crosscheck}
电子全息的相位结果应与下列证据互查：同一区域的 HRTEM 或 HAADF 像（确定形貌与界面位置）、EELS 测得的厚度参数 $t/\lambda$（校验式~\eqref{eq:ch02c-holo} 的薄样品前提）、以及宏观催化性能数据（例如 Au/TiO$_2$ 体系中界面电荷传递的推断是否与低温 CO 氧化的动力学特征一致）。缺少后两者时，「界面处存在电荷传递」只能作为假说提出。
\end{crosscheck}

\subsubsection{原位环境透射电镜与气体池}

多相催化反应是一个动态过程，催化剂与气体反应物分子在原子尺度上相互作用，活性位的原子结构对其催化性能具有决定性作用。因此，催化剂活性位微观结构的直接观测对于研究催化过程的本质至关重要。这正是前几项技术都无法回答的问题：球差校正、三维重构与电子全息提升的都是「静态结构看得多清楚」，而催化关心的是「结构怎么变」。源文记载，Boyes 和 Gai 等分别于 1997 年和 1998 年开始使用原位环境透射电子显微镜（ETEM）实时监测催化反应；研究人员将气相反应在透射电镜腔体内进行，采用微分泵进样，从而避免气体分子进入电子枪的超高真空。

此处已对照原页面图像（书内 p.154 / PDF p.173）：原书 \S 2.7.5 首段原文为「Boyes 和 Gai 等分别于 1997 年和 1998 年开始使用原位环境透射电子显微镜（ETEM）实时监测催化反应」，两位作者开始使用 ETEM 的年份确为 1997 年与 1998 年，与正文一致；分片源文本此处 OCR 只残存「分别于……和 1998 年」，1997 年由原页补足。

\begin{principle}
原位环境电镜与传统透射电镜（见 \S 2.1）在构造上的主要区别是引入了差分抽气系统：在催化剂附近引入气体，并利用微电路 MEMS 技术加热实现催化反应；再在远离催化剂的上下方向通过抽气把这些气体排除至电镜外，从而实现原位观察催化剂在反应中的结构演化，同时维持显微镜关键部位（电子枪、照明系统等）的真空度。具体做法是在物镜、中间镜等处设置抽气通道，借助外置真空泵（主要为分子泵）的连续抽气作用将气体排出镜筒外。

真空系统设计的关键，是利用气体分子从样品附近向远端扩散过程中需要经过多级光阑这一事实：光阑的微米级孔径使得气体分子通过时两侧产生气压差，从而使气体分压自样品向镜筒两端依次递减（即真空度递增），最终确保气体在扩散到电子枪之前被真空泵全部排出，从而保证电镜稳定工作。也就是说，原位环境电镜不是「让整个镜筒充满气体」，而是「在样品附近制造一个局部高压区，再用多级压差把它与镜筒的高真空隔开」——所有压力上限的根源都在这里。
\end{principle}

\begin{measure}
受电镜真空抽气系统能力的限制，同时考虑到气体分子散射电子带来的光束质量下降问题，环境电镜中催化剂附近允许的气体压力存在限制。以 FEI 公司的 Titan 系列环境电镜为例，差分抽气系统即便在使用 10 台以上分子泵的情况下，允许气体压力仍然限制在 20\,mbar；对于小分子气体如 H$_2$、He 等，允许压力往往低于 10\,mbar。

此处已对照原页面图像（书内 p.155 / PDF p.174）：原书原文为「以 FEI 公司的 Titan 系列环境电镜为例，差分抽气系统即便在使用 10 台以上分子泵的情况下，允许气体压力仍然限制在 20\,mbar，对于小分子气体如 H$_2$、He 等，允许压力往往低于 10\,mbar」。分片源文本此处 OCR 把「低于 10\,mbar」残为「往 10 mbar」、并把厂牌型号截去；此前正文据残文写作「以某系列环境电镜为例……允许压力往往为 10\,mbar 量级」，方向与原文相反，现按原页改为「低于 10\,mbar」并补回厂牌型号。

如果高于这一典型值，后果有两个方面：一方面是气体会破坏真空而引发设备停机的风险；另一方面，大量气体分子对电子束的吸收与散射会使得电镜分辨率急剧下降，以至无法实现催化剂显微结构的观察，从而丧失原位显微表征的能力。压力上限因此不是设备参数表的余量，而是「分辨率—压力」权衡线的位置。
\end{measure}

\begin{figure}[htbp]
\centering
\begin{tikzpicture}[x=1mm,y=1mm,font={\fontsize{8}{9.6}\selectfont}]
% 对数压力轴：1 个数量级 = 5.2 mm，p 从 1e-10 到 1e6 mbar
\draw[gray!70] (0,-6) -- (83.2,-6);
\draw[gray!70] (0,-6) -- (0,-7.6);
\node[anchor=north] at (0,-8) {$10^{-10}$};
\draw[gray!70] (10.4,-6) -- (10.4,-7.6);
\node[anchor=north] at (10.4,-8) {$10^{-8}$};
\draw[gray!70] (20.8,-6) -- (20.8,-7.6);
\node[anchor=north] at (20.8,-8) {$10^{-6}$};
\draw[gray!70] (31.2,-6) -- (31.2,-7.6);
\node[anchor=north] at (31.2,-8) {$10^{-4}$};
\draw[gray!70] (41.6,-6) -- (41.6,-7.6);
\node[anchor=north] at (41.6,-8) {$10^{-2}$};
\draw[gray!70] (52,-6) -- (52,-7.6);
\node[anchor=north] at (52,-8) {$1$};
\draw[gray!70] (62.4,-6) -- (62.4,-7.6);
\node[anchor=north] at (62.4,-8) {$10^{2}$};
\draw[gray!70] (72.8,-6) -- (72.8,-7.6);
\node[anchor=north] at (72.8,-8) {$10^{4}$};
\draw[gray!70] (83.2,-6) -- (83.2,-7.6);
\node[anchor=north] at (83.2,-8) {$10^{6}$};
\node[anchor=north west] at (-34,-8) {p/mbar};
% 压力窗口（矩形条）
\fill[gray!25] (0,-2.5) rectangle (20.8,2.5);
\draw[gray!60] (0,-2.5) rectangle (20.8,2.5);
\fill[gray!25] (5.2,6.5) rectangle (20.8,11.5);
\draw[gray!60] (5.2,6.5) rectangle (20.8,11.5);
\fill[gray!25] (36.4,15.5) rectangle (53.6,20.5);
\draw[gray!60] (36.4,15.5) rectangle (53.6,20.5);
\fill[gray!25] (67.6,24.5) rectangle (69.2,29.5);
\draw[gray!60] (67.6,24.5) rectangle (69.2,29.5);
\fill[gray!25] (67.6,33.5) rectangle (79.6,38.5);
\draw[gray!60] (67.6,33.5) rectangle (79.6,38.5);
% 行标签
\node[anchor=east,align=right,text width=32mm] at (-2,0) {模型催化／表面科学 UHV};
\node[anchor=east,align=right,text width=32mm] at (-2,9) {电镜镜筒、电子枪};
\node[anchor=east,align=right,text width=32mm] at (-2,18) {ETEM 样品区};
\node[anchor=east,align=right,text width=32mm] at (-2,27) {封闭气体池（MEMS）};
\node[anchor=east,align=right,text width=32mm] at (-2,36) {工业催化反应};
% 条旁标注
\node[anchor=west,align=left,text width=24mm] at (55,18) {上限 $\approx$ 20\,mbar};
\node[anchor=west,align=left,text width=14mm] at (71,27) {1--2\,atm};
\node[anchor=east,align=right,text width=30mm] at (66,36) {常压--200\,bar};
% 压力落差指示
\draw[dashed,gray!70] (53.6,20.5) -- (53.6,31.5);
\draw[dashed,gray!70] (67.6,29.5) -- (67.6,31.5);
\draw[<->,gray!70] (53.6,31.5) -- (67.6,31.5);
\node[anchor=east,align=right,text width=26mm] at (52,30) {约 2--4 个数量级};
% 轴下说明
\node[anchor=north west,align=left,text width=112mm] at (-34,-14) {镜筒要求的高真空与工业催化反应（常压--200\,bar）之间隔着 9--11 个数量级。差分抽气把 ETEM 样品区提到约 20\,mbar，与工业条件仍差 2--4 个数量级；封闭气体池把可及压力提到 1--2\,atm，压力落差被压缩到 1 个数量级以内，代价是电子束两次穿过氮化硅薄膜。};
\end{tikzpicture}
\caption{透射电镜可及气体压力范围与真实催化反应条件的对照（对数标度，数值为源文给出的典型值而非某一机型的保证指标）。读者应看的不是单条数值，而是三个窗口之间的间距：镜筒真空到工业反应隔着 9--11 个数量级，差分抽气原位环境电镜把样品区提到约 20\,mbar 后仍差 2--4 个数量级，封闭气体池则把可及压力提到 1--2\,atm。压力只是「环境落差」的一个维度，温度与气体组分见 \tabref{tab:ch02c-pt}。\figsrc{据 S1 \S 2.7.5，书内 p.155--157 / PDF p.174--176 重绘}\schemamark}
\label{fig:ch02c-pressure}
\end{figure}

除差分抽气系统外，原位环境电镜还有两项配套设计，它们各自也带来新的限制。

其一是污染控制。考虑到催化反应在样品室内壁引发的分子凝聚、积炭等电镜污染问题，往往会在样品室附近配备导流式远程等离子体清洗装置，借助电镜真空系统抽气气流，将产生的等离子体注入样品室，借助氧化电离作用清除内壁上的污染物，以确保电镜的真空度和显微性能。注意这里已经埋下一条重要边界：反应体系中同时存在等离子体，而等离子体本身也会与样品作用。

其二是气体组分监测。除了在高空间分辨条件下观察催化剂结构变化，监测反应过程的气体组分变化同样是全方位获取催化反应信息不可或缺的一环。环境电镜在样品室附近设计了残余气体分析系统（residual gas analysis，RGA）：在样品附近的抽气管道，借助分子泵将少量催化剂附近的气体采样并送至质谱仪检测，给出反应中的底物和产物种类、含量等信息。值得注意的是，RGA 分析是把气体采样至镜筒外的质谱仪，由于气体跨距离传输会带来信号分析的时延，反应气体分析和催化剂显微表征难以做到完全同步。

与 RGA 对应的一种气体检测方案是借助能量损失谱：利用样品室内气体分子对电子束的能量吸收对应于气体分子轨道跃迁的特性，在 EELS 谱低能区（low loss）可以观察到对应气体分子电子能量吸收现象，这些吸收峰往往成组出现，且随气体分子不同而具有特定强度分布，故称为气体分子的「指纹峰」。通过辨识 EELS 谱中指纹峰的种类及强度，可以实现气体组分的辨识乃至定量。由于该信号是从电子束与样品附近气体直接作用产生的（原位、实时产生和记录），这一方案可以实现催化剂表面反应过程的实时、原位监测。然而它不利的一点同样明确：电镜光路工作在能量损失谱模式而非成像模式，导致分析气体组分时无法实现催化剂显微结构的成像观察。于是，「同一时刻只能看结构或只能看气氛」成为原位环境电镜的第一条硬边界。

在电镜中实现催化剂原位气相反应观察的另一个途径是借助原位气体池技术（gas cell）。该技术基于常规透射电镜，不要求在电镜中直接引入气体环境，而是借助特殊的样品杆前端设计：用 MEMS 技术特制的芯片实现精确控温，两片对电子束透明的氮化硅（Si$_3$N$_4$）薄膜将催化剂粉体封闭在微腔内，微腔设有进气、出气孔道，通过样品台与电镜外控制装置连通，通入特定气体并将其约束在反应微腔内，从而不破坏电镜样品室的真空工作环境，实现在常规电镜下的气氛反应，并透过氮化硅薄膜观察催化剂颗粒的微结构演化。

气体池解除的正是 ETEM 的压力上限：由于它需要构筑上下密封的 Si$_3$N$_4$ 薄膜微腔，薄膜强度决定了通入气体的压力，最大典型压力值可达 1--2 个大气压，优于原位环境电镜的 20\,mbar。但它引入的新限制有三条：一是为保持氮化硅薄膜的必要机械强度而增加厚度，反过来会干扰电子束的成像质量，使电镜分辨率下降；二是与 ETEM 相比，气体池中样品两侧的氮化硅薄膜会干扰 X 射线能谱和电子能量损失谱，影响催化剂的成分及电子态分析；三是基于 MEMS 技术的微腔气密性如果失控，可能带来气体泄漏，为电镜的安全运转带来隐患。

\begin{table}[htbp]
\centering
\setlength{\tabcolsep}{4pt}
\small
\caption{原位气氛反应的压力与温度边界：两种原位方案与工业反应条件的对照。读者应重点比较「压力落差」与「分析通道」两行——压力落差最小的方案（气体池）付出了薄膜干扰谱学分析的代价，而谱学通道最完整的方案（ETEM）压力上限最低。表中数值为源文给出的典型值，小分子气体的压力上限见正文的存疑标注。\src{据 S1 \S 2.7.5，书内 p.155—160 / PDF p.174—179 归纳。}}
\label{tab:ch02c-pt}
\begin{tabular}{@{}p{22mm}p{32mm}p{30mm}p{28mm}p{30mm}@{}}
\toprule
条件项 & 原位环境电镜（ETEM） & 封闭气体池（MEMS） & 工业催化反应 & 对结论的影响 \\
\midrule
反应区压力 & 典型上限约 20\,mbar；小分子气体为 10\,mbar 量级（见正文存疑标注） & 最大典型值 1--2 个大气压 & 常压--200\,bar（$10^{3}$--$10^{5}$\,mbar） & 压力改变表面覆盖度与物相稳定性，低压下的「活性相」未必是工业条件下的活性相 \\
\addlinespace
与工业压力落差 & 2--4 个数量级 & 1 个数量级以内 & --- & 落差越大，原位结论向工业条件外推的风险越高 \\
\addlinespace
观察与分析通道 & 无窗，可直接做 EDS/EELS，但与成像模式互斥 & 电子束两次穿过 Si$_3$N$_4$ 薄膜，干扰 X 射线能谱与 EELS & 不在电镜内 & 气体池测得的成分与价态需扣除薄膜影响后才可与无窗结果比较 \\
\addlinespace
加热与测温 & MEMS 芯片加热；早期芯片温度不均匀性可达近 80\,\degC & 同 MEMS 芯片控温 & 反应器整体控温，测温点在床层或壁面 & 测温误差直接进入构效关系，是原位结论最脆弱的一环 \\
\addlinespace
产物／组分监测 & RGA（跨距离传输有滞后）或 EELS 低能指纹峰 & 无在线产物分析 & 在线色谱／质谱，与反应器同步 & 电镜内的组分信息与结构信息往往不同步 \\
\bottomrule
\end{tabular}
\end{table}

在催化研究的实例方面，目前原位环境透射电镜已被广泛用于监测工作状态下的工业催化剂，其化学气氛包括还原（氢气、烃类、一氧化碳等）、氧化（氧气、二氧化碳等）或水蒸气。源文认为其中对钒磷复合氧化物催化剂的研究尤为成功：所选目标反应是正丁烷选择氧化制马来酐，氧钒基焦磷酸盐相（(VO)$_2$P$_2$O$_7$，简称 VPO）的 (201) 晶面被认为具有最高的催化活性，该活性相主要通过焙烧前驱体 VO(HPO$_4$)$\cdot$0.5H$_2$O（VHPO）制得，因此对该前驱体向活性相转化过程进行原位研究具有重要意义。在原位环境透射电镜中，将 VHPO 前驱体在氮气氛围下加热后发现，前驱体结构转变从 400\,\degC 左右开始，绝大部分的 VPO 在 450\,\degC 下生成；结合动态电子衍射分析发现，在整个结构变化过程中，所有晶体的结晶程度较高，原子排列的周期性良好，未观测到其他过渡晶相或非晶相的生成。在正丁烷气氛和典型的反应温度（约 390\,\degC）下，只有 VPO 的 (201) 晶面上的阳离子对烷烃选择氧化起作用，其微结构介于氧钒八面体和磷酸四面体之间。使用原位环境透射电镜进行催化研究的实例还包括铜催化水蒸气重整反应和纳米碳管生长。

原子分辨尺度的原位观察则更能说明「结构确在变化」。源文记载，可以观察到反应前 CeO$_2$ 负载的 Au 单晶纳米颗粒保持标准的 (100) 面间距 0.20\,nm，同时面内原子距离为 0.29\,nm；而在催化 CO 氧化过程中，表面 CO 分子吸附导致 Au 的 (100) 晶面原子发生重构，面内原子间距被压缩为 0.25\,nm，面间距离则拉伸至 0.25\,nm。此处已对照原页面图像（书内 p.158 / PDF p.177）：原书原文为「可以观察到反应前 Au 单晶纳米颗粒保持了标准的（100）面间距 0.20 nm，同时面内原子距离为 0.29 nm；而在催化 CO 氧化过程中，表面 CO 分子吸附导致 Au 的（100）晶面原子发生重构，面内原子间距被压缩为 0.25 nm，面间距离则拉伸至 0.25 nm」。重构前后两处 0.25\,nm 确为原书原文——图 2-56 中亦分别标注 0.29\,nm／0.20\,nm（真空）与 0.25\,nm／0.25\,nm（CO／空气），并非 OCR 失真；面内值原书记作 0.29\,nm（此前正文误作 0.291\,nm，已改正），「特征间距」原书作「（100）面间距」，「层间距离」原书作「面间距离」。由此可见，ETEM 手段为在原子尺度探测催化剂表面重构行为提供了有力保障。源文进一步指出，在当前技术条件下，如匹配超快电子感光相机（基于 CMOS 的电子感光元件），有望在毫秒乃至微秒尺度捕捉催化剂在反应中的原子级结构变化。另一个例子是追踪 Fe 单个颗粒在费托反应中表面 Fe$_5$C$_2$ 相的演化过程：在环境电镜中引入 CO 与 H$_2$ 的反应气氛，借助单颗粒电子衍射及暗场像技术，在反应 60\,min 内追踪了 Fe 颗粒的表面碳化、Fe$_5$C$_2$ 成型及生长过程。此处已对照原页面图像（书内 p.159 / PDF p.178）：原书原文为「在环境电镜中引入 CO 与 H$_2$ 的反应气氛，借助单颗粒电子衍射及暗场像技术［图 2-57（b）］，在反应 60 min 内追踪了 Fe 颗粒的表面碳化、Fe$_5$C$_2$ 成型及生长过程［图 2-57（a1）～（a3）］」。时间单位为 \textbf{min}（图 2-57（b）中标出的时刻为 8、9、12、40、42、60 min），此前正文因 OCR 只剩「在反应 60～ 内」而写作「60\,s 量级」，与原书不符，已改正并补出「表面碳化、Fe$_5$C$_2$ 成型及生长」三阶段与「单颗粒电子衍射及暗场像技术」。

\begin{obs}
在 CeO$_2$ 负载 Au 颗粒的 CO 氧化原位实验中，可从图像序列直接读出：（1）反应前颗粒呈单晶形貌，面内原子距离为 0.29\,nm（原书图 2-56 标值）；（2）通入 CO／空气并升温后，同一颗粒表面的原子列间距变小（面内压缩），并且沿另一方向的间距发生变化；（3）变化发生在外表面，颗粒内部仍保持原有排列。
\end{obs}

\begin{infer}
上述变化被解释为表面吸附导致 Au 的 (100) 晶面发生原子重构。这条推断依赖两个前提：一是图中所测间距的变化确实来自原子列位置而不是衬度条件的改变（欠焦、厚度、通道效应都会改变表观间距）；二是重构被归因于反应气氛，而不是电子束本身。后一个前提尤其需要对照实验来保证——见下面的误判与回查。
\end{infer}

\begin{dataflow}
原位实验的完整数据链是：设定气氛（气体种类、分压、流速）与温度（设定值 $+$ 实测值 $+$ 不确定度）$\to$ 同步采集图像／衍射／谱（并记录时间戳）$\to$ 结构变化的判读（相变温度、暴露晶面、重构前后间距）$\to$ 与同批样品的产物分析（RGA、EELS 指纹峰或离线色谱）对齐时间轴 $\to$ 建立构效关系假设 $\to$ 用宏观反应器数据检验。链条中最容易被省略、也最致命的一步是「时间轴对齐」：RGA 的传输时延与 EELS 的模式互斥，都会让「结构变化」与「产物变化」看起来同时发生，而实际并不同时。
\end{dataflow}

\begin{pitfall}
源文专列了一段「原位环境电镜表征中的注意事项」，其中每一条都可以理解为这项技术的固有代价：

（1）辐照损伤在高温高压条件下更严重。高温高压反应条件下催化剂晶体结构的电子辐照损伤较为严重；球差校正与三维重构中得到的高分辨图像同样伴随高剂量。判断「这是反应造成的结构变化还是电子束造成的」，必须做束流对照（同一条件、不同剂量率、或束流关断—开启循环）。

（2）电子束会改变气氛。气体分子可能在高速电子束作用下离子化，产生的等离子体与样品发生作用，大大影响了对催化反应的研究。这意味着电镜内的「反应条件」并不是中性的观察窗口，电子束本身是反应体系的一部分。这一条与前述等离子体清洗装置叠加在一起，使样品室内始终存在非热平衡的活性物种。

（3）温度测量的准确性与均匀性。原位环境催化反应需要在特定温度下进行，因此催化剂颗粒需要首先负载到带有 MEMS 控温芯片的氮化硅薄膜上；芯片温度测量的准确与否，直接关系到原位反应中催化剂构效关系的建立，以及与真实反应条件的参考和关联。

源文给出了一种巧妙的原位测温思路：金属受热后原子间距与电荷密度改变，进而引发体相等离子体振动频率移动；反之，借助 EELS 表征等离子体振动频移，就可以反推出 MEMS 加热微电路中的温度。同时源文提醒，对于早期生产的用于 ETEM 的控温芯片，温度均匀性令人担忧，温度差异可达近 80\,\degC，这将导致实际参数的严重偏离，误导对催化剂构效关系的认知；值得庆幸的是，随着 MEMS 技术及工艺的不断进步，控温精度及均匀度在不断提高。80\,\degC 这个数字应当被记住：在 Arrhenius 意义上，它足以把速率与选择性的比较彻底打乱，而它出现在「测温」这一环节，正好说明原位实验的误差不总是来自电镜的分辨率。
\end{pitfall}

\begin{crosscheck}
原位结论必须用三条外部证据回查：一是气氛信息的两条独立通道互查——RGA（镜筒外质谱，有传输时延）与 EELS 低能指纹峰（原位实时，但与成像模式互斥），两者给出的组分变化在时间轴对齐后应当自洽；二是两种原位方案的互查——开腔 ETEM（压力低、无窗、谱学好）与封闭气体池（压力高、有窗、谱学受薄膜干扰）在重叠压力区间内应给出相同的结构演化趋势，若不同则优先怀疑薄膜效应或压力效应；三是与宏观反应器的互查——把电镜条件下的温度、分压换算成与固定床可比的量，检查产物分布与动力学是否一致。三条互查都通过时，原位观察才构成机理证据；否则它只是「在电子束下、在某一气氛中发生过的结构变化」。
\end{crosscheck}

\paragraph{本节小结} 表~\ref{tab:ch02c-newold} 与表~\ref{tab:ch02c-situ} 把本节六项技术的「解除的旧限制／引入的新限制／适用样品」并列在一起。使用时请从右往左读：先确认样品是否落在「适用样品」一列，再确认自己是否具备「前提条件」，最后才看「解除的旧限制」这一列的宣传性表述。

\begin{table}[htbp]
\centering
\setlength{\tabcolsep}{4pt}
\small
\caption{成像与能量分辨类新技术的新旧限制对照。读者应看「引入的新限制」一列：每一项能力的代价都写在那里，而它往往就是实际实验中最先遇到问题的地方。\src{据 S1 \S 2.7.1—\S 2.7.2，书内 p.147—151 / PDF p.166—170 归纳。}}
\label{tab:ch02c-newold}
\begin{tabular}{@{}p{26mm}p{36mm}p{44mm}p{34mm}@{}}
\toprule
技术 & 解除的旧限制 & 引入的新限制 & 适用样品与前提 \\
\midrule
成像球差校正器 & 物镜球差造成的漫散圆斑与菲涅耳条纹；不必再用焦系列计算重建换取清晰度 & 只校正三阶球差，色差与更高阶几何像差成为新的分辨率瓶颈；校正器本身需要稳定的供电与环境；校正后的像仍经中间镜继续放大 & 薄样品、需晶格分辨的纳米晶与界面；束敏样品必须改用低剂量方案 \\
\addlinespace
束斑球差校正器（probe corrector） & STEM 探针无法同时做到最小束斑与足够强度，限制 HAADF 的 $Z$ 衬度分辨率（源文：0.24\,nm 提高到 0.13\,nm） & 需要计算机自动调校多级四极—八极系统；探针电流集中于极小体积，单位面积剂量升高；轻元素柱的衬度仍受 $Z$ 衬度趋势限制 & 原子分辨 HAADF-STEM、原子级 EELS/EDS、低加速电压下的束敏样品 \\
\addlinespace
Wien 单色器与高分辨能量过滤器 & 热场发射系统 0.8--1\,eV 的能量分辨率使 ELNES 精细结构湮没在谱中 & 单色器滤除大量电子（分辨率优于 0.1\,eV 时约 75\% 的场发射电子被滤掉），束流与信噪比下降；低能区定量受相对论效应与切伦科夫辐射干扰 & 过渡金属 $L_{2,3}$ 边、氧 $K$ 边等近阈精细结构；纳米半导体带隙（需修正切伦科夫项） \\
\bottomrule
\end{tabular}
\end{table}

\begin{table}[htbp]
\centering
\setlength{\tabcolsep}{4pt}
\small
\caption{三维与原位类新技术的新旧限制对照。两类技术的共同点是：解除的限制写在中间一列，而新限制往往由「样品台／样品杆」这一类硬件引入，不在电镜主体的指标表上。\src{据 S1 \S 2.7.3—\S 2.7.5，书内 p.151—160 / PDF p.170—179 归纳。}}
\label{tab:ch02c-situ}
\begin{tabular}{@{}p{26mm}p{36mm}p{44mm}p{34mm}@{}}
\toprule
技术 & 解除的旧限制 & 引入的新限制 & 适用样品与前提 \\
\midrule
三维电子断层重构 & 二维投影把三维结构特征相互覆盖，无法判断金属粒子在载体表面还是孔道内部 & 需 $\pm 80^\circ$ 大角度倾转，物镜极靴空间加大导致物镜分辨率下降（常靠球差校正器补偿）；缺失楔形使沿光轴方向失真；需上百张投影，累积剂量与漂移；原子分辨与元素／价态分辨目前不能兼得 & 多孔载体上金属粒子的空间分布与孔内／孔外比例、核壳结构的价态分区、等离基元模式分布；前提是每张投影的强度随质厚单调 \\
\addlinespace
电子全息术（双棱镜） & 一般高分辨像中相位与振幅非线性混合、衬度受相位传递函数修饰；可绕开物镜球差直接取相位 & 需高亮度、高相干性场发射源；需参考全息图与傅里叶重建；只有薄样品与弱相位近似下相位才与投影电势线性相关；相位需解缠，视场受相干宽度限制 & 纳米粒子的投影电势与内部空洞、带电粒子周围静电场、磁畴；催化体系应用仍有限（如 Au/TiO$_2$ 界面电荷传递） \\
\addlinespace
原位环境电镜（差分抽气） & 电镜高真空与催化反应气氛不相容，无法在反应中实时观察 & 样品区压力上限（典型 20\,mbar；小分子气体 10\,mbar 量级）与工业条件仍差 2--4 个数量级；气体散射电子使分辨率下降；气体电离产生等离子体参与反应；RGA 与成像不同步；EELS 指纹峰与成像模式互斥；高温高压下辐照损伤加重 & 还原／氧化／水蒸气气氛下的结构演化：VPO 前驱体相变、Au/CeO$_2$ 表面重构、Fe 费托碳化物生长 \\
\addlinespace
原位气体池（MEMS 封闭腔） & ETEM 的压力上限；可在常规电镜上实现 1--2 个大气压的气氛反应 & 为保持机械强度而加厚的 Si$_3$N$_4$ 薄膜干扰成像，分辨率下降；薄膜干扰 X 射线能谱与 EELS，影响成分及电子态分析；微腔气密性失控会漏气，威胁电镜安全 & 较高压力下的颗粒形貌与物相演化；不适用于需要无窗 EDS/EELS 定量的场合 \\
\bottomrule
\end{tabular}
\end{table}

\subsection{透射电子显微镜的局限性及注意事项}

同任何其他表征仪器一样，透射电子显微镜也有其局限性。源文把局限列为六条（投影、统计代表性、大小粒子、衬度与结构的关系、放大倍数与晶格参数、定量成分分析的精度）；本节保留这六条，并补上原位与束损伤相关的三条。为便于使用，下面把全部局限归并为五类「落差」，每一条落差都给出它在实验中的表现与缓解方式：投影效应（维数落差）、统计代表性（空间落差）、束损伤（剂量落差）、真空与真实反应条件的落差（环境落差）、样品制备改变样品本身（同一性落差）。最后给出结论前必须报告的项。

\begin{methodsum}
本节的用法不是「记住电镜不可靠」，而是「在写结论时逐条自查」。一个可操作的判据是：任何一条由电镜得出的结论，如果它能被归入下列五类落差之一而作者没有给出应对，就应当被降级为「观察到的现象」，而不是「测得的性质」。表~\ref{tab:ch02c-gaps} 给出了这五类落差的对照与报告要求。
\end{methodsum}

\subsubsection{投影效应：一张像不能回答三维问题}

电子显微像是入射电子透过试样后成的二维投影像，因此用一张电镜照片不能简单地解释试样的三维结构。尤其是催化剂比表面积大、多孔样品多，从一张投影照片不能判断出金属粒子是负载在载体表面还是在载体孔道内部。源文举的例子很典型：用纳米碳管作载体制备金属催化剂，透射电镜只能验证制备是否成功，金属粒子是否负载在纳米碳管上；虽然照片可以给出金属粒子大小、分布、形貌的信息，但并不能准确定位金属粒子是在纳米管内部或外部。在普通透射电子显微镜中，确定金属粒子是否分布在碳管内外的唯一办法就是转动样品（如 $\pm 30^\circ$）观察金属粒子相对于碳管壁的位置变化。

需要强调的是，这里的困难不是分辨率不够，而是信息维度不够：投影把沿束方向不同深度的结构叠加到一个像素上，任何「深度」方向的问题（孔内还是孔外、核壳还是混合物、表面还是体相）都不能由单张投影像回答。解决它的技术是电子断层重构，而重构又带来新限制（见 \S 2.7.3）。

\subsubsection{统计代表性：局域信息对上亿粒子}

催化反应中使用的催化剂是大量的。如果是负载型金属粒子催化剂，参加反应的金属粒子上亿。电子显微镜具有高空间分辨率，但给出的是局域信息：一张高分辨电子显微像通常只包括一两个，最多四五个金属粒子的结构像；即使低放大倍数的显微像照出的粒子数目也经常少于一百个（这与粒子的大小有关：粒子太小时，放大倍数太低不利于观察其形貌并测量其直径）。因此做电子显微镜实验时不能照一张高分辨像就做出任何结论；所照的照片一定要有代表性、重复性，不能选择照「漂亮」的区域。确定催化剂结构或相时，不要忽略利用电子衍射技术：电子衍射花样的强度通常来自于几十个或几百个粒子，这使确定催化剂结构更准确些，并更具代表性。

用电子显微技术确定催化剂的粒子大小分布时，还存在大粒子与小粒子的问题。因为在电子显微照片中大粒子的衬度清晰可辨，统计分布时一般不会漏掉；但是当载体厚度不均时，小于 2\,nm 的粒子衬度低、强度小，不易观察到。因此做统计分布时除了照大量的电子照片（包括至少几百个粒子到上千个粒子）外，还应注意不要忽略统计小粒子的尺寸，但同时也不能主观地忽略不统计大粒子。实验上，如果要把负载粒子的大小分布与催化剂制备和其催化性能联系起来，所做的统计分布必须基于大量粒子，并不能带有主观性。

注意这两条的联合后果：粒径分布是对数正态型的宽分布时，「不漏大粒子、不漏小粒子」在同一个放大倍数下无法同时满足，所以正确的做法是多个放大倍数分层统计，并在报告中给出各层统计的粒子数与判据。

\subsubsection{束损伤：观察本身会改变样品}

除样品的质厚衬度及扫描透射电子显微镜中的 $Z$ 衬度像的强度与试样的厚度或原子序数平方成正比外，电子显微像的衬度与样品的二维投影结构没有直接线性关系。这一条本身是「读图」的限制（见下一条），而与之并行的是「电子束对样品的改写」：源文在讨论原位环境电镜时明确指出，高温高压反应条件下催化剂晶体结构的电子辐照损伤较为严重，并且气体分子可能在高速电子束作用下离子化、产生的等离子体与样品发生作用，大大影响了对催化反应的研究。

以下三点是电子显微学的通行结论，源文未逐条展开，写在这里是为了让读者把「剂量」当成一个必须报告的量（而非可调的经验参数）：（1）入射电子能量（通常 80--300\,keV）远高于化学键能（几个 eV），能量通过弹性与非弹性散射沉积在样品中，可造成原子位移（knock-on）、溅射、局部升温、键断裂、物种迁移与非晶化；（2）氧化物、分子筛、金属有机框架、碳材料与含有机配体的体系对剂量尤其敏感，而催化剂往往正属于这些类别；（3）因此在「看得见」与「没改变」之间存在一条权衡线，只能通过低加速电压、低剂量率、直接电子探测器、以及足够的重复实验来换取，不能通过事后解释来消除。

操作上的最低要求是设置束流对照：在同一区域连续采集若干张像，比较第一张与最后一张；或在同一条件下改变剂量率重复实验。若结构变化只在第二张之后出现，它更可能是束损伤而不是本征结构。

\subsubsection{真空与真实反应条件的落差}

这是原位技术带来的落差，也是把电镜结论外推到催化反应时最容易断裂的一环。镜筒需要高真空（电子枪区在 $10^{-6}$\,mbar 以下量级），而工业催化反应在常压到 200\,bar（$10^{3}$--$10^{5}$\,mbar）进行，两者相差 9--11 个数量级。差分抽气把样品区压力提到约 20\,mbar（图~\ref{fig:ch02c-pressure}），与工业条件仍差 2--4 个数量级；封闭气体池把可及压力提到 1--2\,atm，压力接近了，但样品被两片 Si$_3$N$_4$ 薄膜封闭，电子束与 X 射线都要穿过薄膜，谱学分析的通道被削弱。

更重要的是，压力不是唯一的变量。工业反应器中还存在气体组成梯度、停留时间分布、传热与颗粒团聚状态，以及成型体的尺寸效应；这些在电镜里都不存在。电镜看到的是「单个或少量颗粒、在准静态气氛中、在电子束照射下」的行为。因此原位结果的正确用法是作为机理假设的产生器与证伪工具（例如确认某一活性相在反应气氛下是否稳定），而不是反应器行为的直接测量。凡是用原位像直接给出反应速率、选择性的做法，都需要与宏观反应器数据交叉验证。

\subsubsection{样品制备改变样品本身}

电子透明性要求把样品做成薄区（高分辨观察通常需要几十纳米甚至更薄），而催化剂在使用状态下是微米到毫米级的颗粒或成型体。把样品变成「电镜可见」的过程本身会改变它：研磨与离子减薄会引入缺陷、非晶层与表面终止的改变；超声分散与滴涂会破碎孔结构、改变颗粒的团聚状态与暴露晶面的比例；载体的孔道与金属粒子的空间关系也可能在制样中被改变——这恰好与投影效应耦合在一起，使「粒子在孔内还是孔外」的判断同时受制于制样与投影。

原位实验还有额外一层：为了用 MEMS 芯片控温，催化剂粉体必须负载到氮化硅薄膜或芯片上，其传热方式、气氛接触方式与工业成型体不同。

这一条源文没有单列，本节据 \S 2.7.5 的制样要求与 \S 2.8 中「样品不很薄时分析空间分辨率要下降」的论述归纳而来。它的操作含义是：报告必须写明制样方法（研磨、离子减薄、分散介质、支持膜）、可能的损伤形式，并尽可能做「不同制样方法给出同一关键结论」的对照；把电镜结论与 XRD、物理吸附、化学吸附、反应动力学等其他表征交叉验证，正是为了对冲这一落差。

\subsubsection{衬度与结构之间没有简单的线性关系}

除质厚衬度及 HAADF 的 $Z$ 衬度像的强度与试样厚度或原子序数平方成正比外，电子显微像的衬度与样品的二维投影结构没有直接线性关系。衍射衬度常带有动力学效应（即衍射束的强度会受样品厚度及试样偏离布拉格条件的大小影响），相位衬度受透镜畸变与欠焦量控制。因此高分辨像中的暗（或明）条纹不是试样的某一晶面的原子排列，而只是具有这一晶面间距的条纹像。除非在 Scherzer 欠焦时衬度传递函数在很大的空间频率范围内不改变符号，否则不能把高分辨像的亮点（或暗点）用原子排列位置解释。

这条限制直接决定了读图语言：看到周期性的条纹，可以推断「存在具有该间距的晶面族」，但不能直接推断「亮点的位置就是原子柱的位置」。把条纹像当原子像来标注原子坐标，是本卷最需要避免的过度解释之一。

\subsubsection{放大倍数、晶格参数与定量分析的精度}

电子显微镜的放大倍数（和做电子衍射时的相机长度）只是电子显微镜出厂时厂家给出的参考值。如要用电子显微镜精确测量晶格参数，一定要用标样标定放大倍数。注意在高放大倍数时，透镜激磁电流的任何微小变化都会导致电子显微镜的放大倍数（或测定的晶格间距）变化，因此照高分辨像的实验条件（如室温或冷却水的温度）都应与标定时的实验条件一致。最好的办法是用样品中已知的结构定标来测定晶格常数。另外，小粒子在倒易空间中的格子点变大（尺寸效应），所以在电子显微镜中即使小粒子不完全处于布拉格条件时也能看见其点阵条纹（对面心立方金属最常见的是 $\{111\}$ 晶面的干涉条纹）。但当粒子小于 2\,nm 时，测得的点阵间距要比金属晶体的点阵间距大得多——此时用条纹间距反推物相会系统性出错。

定量成分分析方面，现代的分析电子显微镜的定量成分分析技术已达到小于 1\% 的精确度，但这样的精确度一般都是在研究 X 射线能谱和电子能量损失谱探测极限的专门实验中得到的。因为样品几何条件的限制、计算电离截面时理论值的近似性、切除背底强度引入的误差，在平时表征催化剂试样的实验中通常达不到这样的高度：普通的分析电子显微镜的定量分析若有 5\% 的误差是不奇怪的，而且不能在试样上只取一个谱测定其成分，必须针对不同区域做定量分析并给出误差范围。

空间分辨率一侧同样有落差。现代高空间分辨率的分析电子显微术都能把电子束形成小至 1\,nm 的电子束斑和电子探针来做微区成分定量分析，包括扫描测定试样某一方向直线的化学分析；但是应该注意，当样品不是非常薄时，1\,nm 直径的电子束射到样品上与试样相互作用时，相互作用的体积空间往往要远大于 1\,nm 直径。即由于多重散射，电子束在样品内变宽，所收集的 X 射线光子信号等会来源于几纳米的作用范围，所以真正分析的试样区域通常要比电子束斑大得多。故样品不很薄时，分析电子显微镜的空间分辨率要下降。另外因为入射电子的能量通常在 100\,kV 以上，所以 X 射线能谱测定成分是样品的体相成分，不是其表面相成分——这一点对表面敏感的催化问题尤其重要：由 EDS 得到的「表面组成」其实是体相组成。

\subsubsection{结论前必须报告的项}

\begin{table}[htbp]
\centering
\setlength{\tabcolsep}{4pt}
\small
\caption{透射电镜结论的五类落差及其报告要求。读者应把右列当作投稿或内部报告前的检查清单：任何一项缺失，都要在正文中用「在本实验条件下」之类的限定语把结论收窄。\src{据 S1 \S 2.8，书内 p.160—162 / PDF p.179—181 归纳；其中「样品同一性」一条兼采 \S 2.7.5 的制样要求。}}
\label{tab:ch02c-gaps}
\begin{tabular}{@{}p{20mm}p{28mm}p{30mm}p{30mm}p{28mm}@{}}
\toprule
落差类型 & 电镜一侧 & 真实催化一侧 & 主要后果 & 报告与缓解要求 \\
\midrule
投影（维数） & 单张二维投影像 & 三维多孔颗粒与成型体 & 无法判断粒子在孔内或孔外；核壳与混合物无法区分 & 报告倾角或倾转实验；三维结论给出倾角范围、缺失楔形与体素分辨率 \\
\addlinespace
统计代表性（空间） & 一张像 1--5 个粒子；低倍像常少于 100 个粒子 & 参与反应的金属粒子上亿 & 结论依赖所选区域；$<$2\,nm 粒子易被漏计 & 分层多放大倍数统计、给出粒子总数与判据、禁止只选「漂亮」区域；结构判定用电子衍射交叉 \\
\addlinespace
束损伤（剂量） & 高分辨与谱学都需要高剂量 & 使用状态下的催化剂不承受电子束 & 观察到的变化可能是束诱导的重构、非晶化或物种迁移 & 报告剂量与剂量率；做连续采集与变剂量率对照；束敏体系改用低电压／低剂量方案 \\
\addlinespace
真空与压力（环境） & 镜筒高真空；ETEM 约 20\,mbar；气体池 1--2\,atm & 常压--200\,bar，且有组成梯度与传热 & 低压下的活性相未必是工业条件下的活性相；原位像不能直接给速率与选择性 & 报告压力、温度（含测温方法与不确定度）、气体组成与流速；与宏观反应器数据交叉验证 \\
\addlinespace
样品同一性（制样） & 薄区、分散在支持膜或芯片上 & 微米--毫米级颗粒或成型体 & 表面终止、孔结构、团聚状态与暴露晶面比例被改变 & 报告制样方法（研磨、离子减薄、分散介质、支持膜）与可能损伤；用不同制样方法做关键结论对照；与 XRD、吸附、动力学数据互查 \\
\bottomrule
\end{tabular}
\end{table}

\begin{crosscheck}
把本节压缩成一份可执行的检查清单：结论是原子的还是条纹的（衬度是否允许这样解释）？区域是几个粒子还是上千个粒子（统计是否成立）？剂量与剂量率是多少（是否做过束流对照）？压力与温度离真实条件差多远（含测温不确定度与温度均匀性）？样品是不是原来那个样品（制样与负载是否改变它）？这五问中的每一问，都必须能在报告里找到对应的数字；找不到的，结论就应当加限定语。
\end{crosscheck}

\subsubsection{结束语}

本章介绍了透射电子显微镜的基本构造、操作模式、衬度形成原理及其附属的 X 射线能谱和电子能量损失谱的基本工作原理，并对与电子显微技术有关的最新发展做了一些介绍；用所举的应用实例，讲述了如何应用透射电子显微镜对不同催化剂进行表征，并利用其所具有的高空间分辨率和高能量分辨率研究催化材料的表面结构、缺陷化学及电子能带结构。源文指出，如果能把 TEM、HRTEM、HAADF、EDS 和 EELS 结合起来用于催化材料的研究，会得到比前面讲述的几个例子更多的信息，从而指导催化剂制备、探索催化剂反应机理。电子显微学已发展成为一门专门的学科，要真正掌握这一技术并应用到催化化学中，不仅操作上要熟练，而且还要掌握好晶体学、结构化学、固体物理学和量子力学的知识。随着场发射电子显微镜的日益普及，球差校正器、电子全息术、三维电子显微术及原位环境电镜的发展，电子显微术会成为研究催化化学必不可少的工具。

\srcnote{\S 2.7--2.8 据原书第 2 章，书内 p.147--162 / PDF p.166--181；本节新绘的压力对照图与四张对照表为便于比较而作，图内数值取自源文给出的典型值。正文中凡按学科标准形式重建之处，均已逐处对照原页面图像（书内 p.147--162 / PDF p.166--181）核验，核验结论就地以「此处已对照原页面图像（书内 p.X / PDF p.Y）」的说明记于该处；未加此类说明之处，不代表已核对页面图像。}
